\section{Life of a platform: software updates}
\label{sec:lifecycle}

\begin{inbrief}
A vehicle platform is built for years, while its software keeps growing.
The lifecycle analysis sizes the hardware for the software at the start of
production, then adds the task-chains of successive software updates to
that fixed hardware. It measures how many updates the platform absorbs, and
how its timing margin and deadline satisfaction evolve, against the cost of
the hardware.
\end{inbrief}

\paragraph{The question.} The hardware of a vehicle platform is chosen
before the SOP and is then fixed for the life of the vehicle line. Software
updates, over the air or at service, add functions to it for years. A
platform sized tightly for the SOP software is cheap, but the first update
may not fit. A platform with room to spare costs more cores from the first
car on. The lifecycle variant of the problem (\cref{sec:problem}) makes
this trade-off measurable. It is related to the extensibility of a
distributed system \cite{zhu2009optimizing}, but here the room is consumed
by whole new task-chains, with their own pinned drivers and network
streams, not by growing WCETs.

\paragraph{Procedure.} The SOP mapping is the design point of
\cref{sec:method} for the initial application $\Chains_0$, at a
\emph{provisioning} cap $\varphi_{\mathrm{SOP}}$. A lower cap buys more
cores. The platform is then fixed: the cores this mapping uses, nothing
more. Update $u=1,2,\dots$ brings a set $\Chains_u$ of new task-chains. It
is processed as follows (\cref{alg:lifecycle}).
\begin{enumerate}
\item The new task-chains are partitioned like the initial ones
(\cref{alg:partition}).
\item Their sub-chains are allocated with the rule of the SOP mapping, onto
the cores of the platform only (\cref{alg:assign} with fixed hardware: no
core may be added). The sub-chains already deployed stay where they are,
because an update re-flashes software and should not move what has been
validated. The cores may now be filled up to the cap of the updates,
$\varphi_{\mathrm{upd}}\ge\varphi_{\mathrm{SOP}}$.
\item The priorities of every core are assigned again
(\cref{alg:audsley}). The tasks keep their cores, but their priorities may
change, which a real update would have to re-validate.
\item The update is \emph{absorbed} if every new sub-chain found a core and
every core is schedulable. The end-to-end timing of all task-chains, old
and new, is then analysed again: a new stream can delay old ones on a
shared port.
\end{enumerate}
The first update that is not absorbed ends the life of the platform: its
software would need new hardware. The results are the number of updates
absorbed, and the critical scaling factor, the load and the share of
task-chains that meet their deadline after each update.

\begin{algorithm2e}[t]
\caption{\textsf{Lifecycle}$(\Chains_0,\Chains_1,\dots,\Chains_K)$}
\label{alg:lifecycle}
$\pi\gets$ design point of $\Chains_0$ at $\varphi_{\mathrm{SOP}}$\;
$\Cores\gets$ the cores $\pi$ uses \tcp*{the platform, fixed}
\For{$u\gets1$ \KwTo $K$}{
	$\Frags_u\gets$ \textsf{Partition} of $\Chains_u$\;
	$\pi'\gets\pi$ extended by \textsf{Allocate}$(\Frags_u)$ on $\Cores$, cap $\varphi_{\mathrm{upd}}$, no new core\;
	\lIf{infeasible}{\Return $u-1$ updates absorbed}
	\textsf{Priorities} on every core of $\pi'$\;
	\lIf{a core is unschedulable}{\Return $u-1$ updates absorbed}
	$\pi\gets\pi'$; record $\lambda(\pi)$, load, task-chains on time\;
}
\Return $K$ updates absorbed\;
\end{algorithm2e}

\paragraph{Choices.} The analysis evaluates one update policy, not the
best possible one: a policy that re-allocated deployed software would
absorb more updates (\cref{sec:discussion}). The update does not use consolidation,
relocation moves or the local search, since these would move deployed
software; the deadline satisfaction after an update is therefore that of
the new allocation as it is. The cap of the updates is the design cap of
the platform ($\varphi_{\mathrm{upd}}=0.8$ in the experiments), while the
SOP cap is the provisioning decision. The new task-chains are generated
with the same statistics as the initial ones and numbered after them, so
that their pinned drivers land in random zones, as new functions with new
sensors would.
